\chapter{Category}

\begin{exercise}{1}
    If $f$ is increasing, show that $\omega_f(a) = f(a+) - f(a-)$.
\end{exercise}

\begin{solution}

    First of all, we note that due to $f$ being increasing, $f(a+), f(a-)$ are well-defined and finite: if any of these limits was infinite, we would not be able to define a value for $f(a)$ that would make $f$ increasing.
    For the same reason, $\omega_f(a)$ is also finite.
    Furthermore, because $f$ is increasing, we that, for any $h > 0$, and any $x \in (a - h, a+ h)$

    \[f(a - h) \leq f(x) \leq f(a + h)\]
    
    It must therefore also hold that 
    
    \[\sup_{x, y \in (a - h, a + h)} \lvert f(x) - f(y) \rvert \leq f(a + h) - f(a - h) \implies \omega_f(a) \leq f(a+) - f(a-),\]

    where we have taken limits on both sides.
    Now suppose that $(f(a+) - f(a-)) - \omega_f(a) = \epsilon > 0$.
    Expanding the limits, we obtain:

    \[\lim_{h \rightarrow 0+} f(a + h) - \lim_{h \rightarrow 0+} f(a - h) - \lim_{h \rightarrow 0+} \text{diam} f((a - h, a + h)) = \epsilon > 0 \implies\]

    \[\lim_{h \rightarrow 0+}(f(a + h) - f(a - h) - \text{diam}f((a - h, a + h))) = \epsilon > 0\]

    This means that there exists $\delta > 0$ such that whenever $0 < h < \delta$, it holds that

    \[\lvert f(a + h) - f(a - h) - \text{diam}f((a - h, a+ h)) - \epsilon \rvert < \epsilon \implies\]
    \[f(a + h) - f(a - h) - \text{diam}f((a - h, a + h)) - \epsilon > -\epsilon \implies f(a + h) - f(a - h) > \text{diam}f((a - h, a + h))\]
    \[\implies \sup_{x, y \in (a - h, a + h)} \lvert f(x) - f(y) \rvert < f(a + h) - f(a - h)\]

    Notice then that it must be the case that \textit{at least} one of $f((a + h)-), f((a - h)+)$ does not equal $f(a + h), f(a - h)$ respectively.
    If this was not the case, the supremum would \textit{equal} the RHS.
    We conclude that $f$ has at least one discontinuity for every $h$ in $(0, \delta)$, and so it has uncountably many discontinuities in $(0, \delta$), which, by Theorem 2.17 is a contradiction.
\end{solution}

\newpage

\begin{exercise}{2}
    Prove that $f$ is continuous at $a$ iff $\omega_f(a) = 0$.
\end{exercise}

\begin{solution}
    
    Suppose first that $f$ is continuous at $a$.
    Then, for any $\epsilon > 0$, there exists $\delta > 0$ such that whenever $\lvert x - a \rvert < \delta, \lvert f(x) - f(a) \rvert < \epsilon$.
    We thus have that for any two $x, y \in B_{\delta}(a), \lvert f(x) - f(y) \rvert \leq \lvert f(x) - f(a) \rvert + \lvert f(y) - f(a) \rvert < 2\epsilon$.
    Therefore, $\sup_{x, y \in B_{\delta}(a)} \lvert f(x) - f(y) \rvert \leq 2 \epsilon \implies \text{diam} f(B_{\delta}(a)) \leq 2\epsilon$.
    In other words, the quantity $\text{diam} f(B_{\delta}(a))$ tends to zero.
    Therefore, $\omega_f(a) = \lim_{h \rightarrow 0+} \text{diam} f(B_h(a)) = 0$.
    Conversely, suppose $\omega_f(a) = \lim_{h \rightarrow 0+} \text{diam} f(B_h(a)) = 0$.
    Pick any $\epsilon > 0$.
    Then there exists $\delta > 0$ such that whenever $h < \delta, \text{diam} f(B_{h}(a)) < \epsilon$.
    In other words,

    \[\sup_{x, y \in B_h(a)} \lvert f(x) - f(y) \rvert < \epsilon \implies \sup_{x \in B_h(a)} \lvert f(x) - f(a) \rvert < \epsilon\]

    This means that for all $x$ with $\lvert x - a \rvert < h < \delta, \lvert f(x) - f(a) \rvert < \epsilon$, which means that $f$ is continuous at $a$.
\end{solution}

\begin{exercise}{9}
    If $E$ is a closed set in $\mathbb{R}$, show that $E = D(f)$ for some bounded function $f$.
    [Hint: A sum of two characteristic functions will do the trick.]
\end{exercise}

\begin{solution}
    
    Consider first the functions $g(x) = 0, x \in \mathbb{R}$ and $h(x) = 1$ for $x \in \mathbb{Q}$, $h(x) = -1$ for $x \in \mathbb{R} \setminus \mathbb{Q}$.
    Notice first that $g$ is continuous everywhere, whereas $f$ is discontinuous everywhere (any real can be approached by a sequence of rationals and vice versa).
    We define the following function $f: \mathbb{R} \rightarrow \mathbb{R}$

    \[f(x) = \begin{cases}
        g(x), & x \in \mathbb{R} \setminus E \\
        h(x), & x \in E
    \end{cases}\]

    Notice that $f$ is clearly bounded.
    We claim that $f$ is continuous on all of $\mathbb{R} \setminus E$ and discontinuous on all of $E$.
    To begin, observe that $\mathbb{R} \setminus E$ is an open set.
    Therefore, for any $x \in \mathbb{R} \setminus E$ and any sequence $(x_n), x_n \rightarrow x$, we know that all but finitely many terms of $(x_n)$ are inside $\mathbb{R} \setminus E$.
    Consequently, $(f(x_n))$ eventually becomes the zero sequence, which converges to $f(x) = 0$.
    Therefore $f$ is indeed continuous on all of $\mathbb{R} \setminus E$.
    Now let $x \in E$.
    If $x \in \mathbb{R} \setminus \mathbb{Q}$, then form a sequence $x_n \rightarrow x$ consisting entirely of rationals, which is always possible.
    We then have that $f(x_n) = 0$ or $f(x_n) = 1$ for all $x_n$, and therefore it cannot be that $f(x_n) \rightarrow f(x) = -1$.
    Similarly, suppose $x \in \mathbb{Q}$, and approach $x$ with a sequence $(x_n) \subset \mathbb{R} \setminus \mathbb{Q}$, in which case $f(x_n) = 0$ or $f(x_n) = -1$ and $f(x) = 1$, thus making convergence impossible.
    Therefore, $f$ is discontinuous precisely at every point of $E$, and $D(f) = E$.
\end{solution}
